\setcounter{section}{-1}

\section{Preliminaries}

\subsection{Basics}

For Exercises 1 to 4, let
$$M = 
\begin{pmatrix}
    1 & 1 \\
    0 & 1
\end{pmatrix}$$
and $\bb = \{X \in \ac : MX = XM\}$, where $\ac$ denotes the set of $2 \times 2$ matrices with real entries.
\begin{problems}
    \item Determine which of the following elements of $\ac$ lie in $\bb$:
    $$
    \begin{pmatrix}
        1 & 1 \\
        0 & 1
    \end{pmatrix},
    \begin{pmatrix}
        1 & 1 \\
        1 & 1
    \end{pmatrix},
    \begin{pmatrix}
        0 & 0 \\
        0 & 0
    \end{pmatrix},
    \begin{pmatrix}
        1 & 1 \\
        1 & 0
    \end{pmatrix},
    \begin{pmatrix}
        1 & 0 \\
        0 & 1
    \end{pmatrix},
    \begin{pmatrix}
        0 & 1 \\
        1 & 0 
    \end{pmatrix}$$
    \begin{sol}
        Note that the 1st matrix is $M$ itself so that it belongs to $\bb$, since $MX = XM = M^2$. Further note that the 3rd, 5th, and 6th matrices are the zero, identity, and exchange matrices respectively. Then the 3rd and 5th belong to $\bb$ while the 6th does not. Then only the remaining matrices to check are the 2nd and 4th matrices. For the 2nd matrix:
        \begin{align*}
            MX & = 
            \begin{pmatrix}
                1 & 1 \\
                0 & 1
            \end{pmatrix}
            \begin{pmatrix}
                1 & 1 \\
                1 & 1
            \end{pmatrix} = 
            \begin{pmatrix}
                1 \cdot 1 + 1 \cdot 1 & 1 \cdot 1 + 1 \cdot 1 \\
                0 \cdot 1 + 1 \cdot 1 & 0 \cdot 1 + 1 \cdot 1
            \end{pmatrix} = 
            \begin{pmatrix}
                2 & 2 \\ 1 & 1
            \end{pmatrix} \\
            XM & = 
            \begin{pmatrix}
                1 & 1 \\
                1 & 1 
            \end{pmatrix}
            \begin{pmatrix}
                1 & 1 \\
                0 & 1
            \end{pmatrix} = 
            \begin{pmatrix}
                1 \cdot 1 + 1 \cdot 0 & 1 \cdot 1 + 1 \cdot 1 \\
                1 \cdot 1 + 1 \cdot 0 & 1 \cdot 1 + 1 \cdot 1 
            \end{pmatrix} =
            \begin{pmatrix}
                2 & 1 \\
                2 & 1
            \end{pmatrix}
        \end{align*}
        Then $MX \neq XM$ and the 2nd matrix does not belong to $\bb$. For the 4th matrix:
        \begin{align*}
            MX & = 
            \begin{pmatrix}
                1 & 1 \\
                0 & 1 
            \end{pmatrix}
            \begin{pmatrix}
                1 & 1 \\
                1 & 0
            \end{pmatrix} =
            \begin{pmatrix}
                1 \cdot 1 + 1 \cdot 1 & 1 \cdot 1 + 1 \cdot 0 \\
                0 \cdot 1 + 1 \cdot 1 & 0 \cdot 1 + 1 \cdot 0
            \end{pmatrix} =
            \begin{pmatrix}
                2 & 1 \\
                1 & 0
            \end{pmatrix} \\
            XM & = 
            \begin{pmatrix}
                1 & 1 \\
                1 & 0
            \end{pmatrix}
            \begin{pmatrix}
                1 & 1 \\
                0 & 1
            \end{pmatrix} = 
            \begin{pmatrix}
                1 \cdot 1 + 1 \cdot 0 & 1 \cdot 1 + 1 \cdot 1 \\
                1 \cdot 1 + 0 \cdot 0 & 1 \cdot 1 + 0 \cdot 1
            \end{pmatrix} = 
            \begin{pmatrix}
                1 & 2 \\
                1 & 1
            \end{pmatrix}
        \end{align*}
        So that $MX \neq XM$.
    \end{sol}
    \item Prove that if $P, Q \in \bb$, then $P + Q \in \bb$ where $+$ denotes the usual sum of two matrices.
    \begin{sol}
        For Exercises 2 and 3, let
        \[P = 
        \begin{pmatrix}
            a & b \\
            c & d
        \end{pmatrix}, \quad Q = 
        \begin{pmatrix}
            e & f \\
            g & h
        \end{pmatrix}, \quad P + Q = 
        \begin{pmatrix}
            a + e & b + f \\
            c + g & d + h
        \end{pmatrix}, \quad P \cdot Q = 
        \begin{pmatrix}
            ae + bg & af + bh \\
            ce + dg & cf + dh
        \end{pmatrix}\]
        where $PM = MP$ and $QM = MQ$. Then
        \begin{align*}
            M(P + Q) & = 
            \begin{pmatrix}
                1 & 1 \\
                0 & 1
            \end{pmatrix}
            \begin{pmatrix}
                a + e & b + f \\
                c + g & d + h
            \end{pmatrix} \\
            & = 
            \begin{pmatrix}
                a + e + c + g & b + f + d + h \\
                c + g & d + h
            \end{pmatrix} \\
            & = 
            \begin{pmatrix}
                a + c & b + d \\
                c & d
            \end{pmatrix} + 
            \begin{pmatrix}
                e + g & f + h \\
                g & h
            \end{pmatrix} \\
            & = MP + MQ \\
            & = PM + QM \\
            & = 
            \begin{pmatrix}
                a & a + b \\
                c & c + d
            \end{pmatrix} + 
            \begin{pmatrix}
                e & e + f \\
                g & g + h
            \end{pmatrix} \\
            & = 
            \begin{pmatrix}
                a + e & a + e + b + f \\
                c + g & c + g + d + h
            \end{pmatrix} \\
            & = 
            \begin{pmatrix}
                a + e & b + f \\
                c + g & d + h
            \end{pmatrix}
            \begin{pmatrix}
                1 & 1 \\
                0 & 1
            \end{pmatrix} \\
            & = (P + Q)M \qh
        \end{align*}
    \end{sol}
    \item Prove that if $P, Q \in \bb$, then $P\cdot Q \in \bb$ where $\cdot$ denotes the usual product of two matrices.
    \begin{sol}
        Using $PQ$, proceed as we did above with $M(PQ)$. Rewriting the entries will result in $(MP)Q$. Since $P \in \bb$, then we have $(PM)Q$. We rewrite entries again to result in $P(MQ)$. Because $Q \in \bb$, we have $P(QM)$, and a final rewrite results in $(PQ)M$.
    \end{sol}
    \item Find conditions on $p, q, r, s$ which determine precisely when $
    \begin{pmatrix}
        p & q \\
        r & s
    \end{pmatrix} \in \bb$.
    \begin{sol}
        Let $X$ be the matrix described above. Note that
        \begin{align*}
            MX & =
            \begin{pmatrix}
                1 & 1 \\
                0 & 1
            \end{pmatrix}
            \begin{pmatrix}
                p & q \\
                r & s
            \end{pmatrix} =
            \begin{pmatrix}
                p + r & q + s \\
                r & s
            \end{pmatrix} \\
            XM & = 
            \begin{pmatrix}
                p & q \\
                r & s
            \end{pmatrix}
            \begin{pmatrix}
                1 & 1 \\
                0 & 1
            \end{pmatrix} = 
            \begin{pmatrix}
                p & p + q \\
                r & r + s
            \end{pmatrix}
        \end{align*}
        Because $X \in \bb$, then we may compare entries to obtain the following:
        \[
        \begin{cases}
            p + r = p \\
            q + s = p + q \\
            r = r \\
            s = r + s
        \end{cases}
        \]
        The first and fourth equations force $r = 0$, and the second equation forces $p = s$. Then $\bb$ is classified as
        \[\bb = \set*{\left.
        \begin{pmatrix}
            p & p + q \\
            0 & p
        \end{pmatrix}\,\right|\, p, q \in \r} \qh\]
    \end{sol}
    \item Determine whether the following functions $f$ are well defined:
    \begin{problems}
        \item $f : \q \to \z$ defined by $f(a/b) = a$.
        \item $f : \q \to \q$ defined by $f(a/b) = a^2/b^2$.
    \end{problems}
    \begin{solalph}
        \item No, because $1/2 = 2/4$, but $f(1/2) = 1$ and $f(2/4) = 2$.
        \item Yes; suppose $a/b = c/d$. Then $ad = bc$, or that $a^2d^2 = b^2c^2$. Then $a^2/b^2 = c^2/d^2$, or $f(a/b) = f(c/d)$. 
    \end{solalph}
    \item Determine whether the function $f: R^+ \to \z$ defined by mapping a real number $r$ to the first digit to the right of the decimal point in a decimal expansion of $r$ is well defined.
    \begin{sol}
        No. Note that $1=1.000\ldots = 0.999\ldots$, but $f(1.000\ldots) = 1$ and $f(0.999\ldots) = 9$.
    \end{sol}
    \item Let $f: A \to B$ be a surjective map of sets. Prove that the relation
    $$a \sim b \iff f(a) = f(b)$$
    is an equivalence relation whose equivalence classes are the fibers of $f$.
    \begin{sol}
        The relation above is an equivalence relation, since $=$ is already an equivalence relation on $B$.
        
        Consider now some $b \in B$. Since $f$ is surjective, there exists $a \in A$ such that $f(a) = b$. The equivalence class of $a$ is the set $\set{x \in A \mid x \sim a}$. By definition of $\sim$, this is equal to the set $\set{x \in A \mid f(x) = f(a) = b}$, which is precisely the fiber of $f$ over $b$.
    \end{sol}
\end{problems}